\chapter{Lo \textit{stage} come innovazione}
\label{chap:strategia}

\textit{Nel presente capitolo approfondisco le motivazioni strategiche alla base della proposta di tirocinio, inserendola nel più ampio quadro della visione aziendale di Ergon Informatica. Mediante l'analisi degli obiettivi formativi e dei vincoli progettuali, analizzo e illustro come il tirocinio si inserisca nella strategia aziendale di innovazione e di miglioramento continuo dei processi, configurandosi sia come occasione di crescita per il candidato sia come strumento attraverso cui l'organizzazione esplora nuove soluzioni tecnologiche. Traggo le considerazioni riportate principalmente dalle interviste svolte con il tutor aziendale e dalle osservazioni raccolte durante il periodo di permanenza in azienda.}

\section{Il valore dello \textit{stage} in Ergon Informatica}
\label{sec:valore_stage}

Come discusso nella Sezione~\ref{subsec:ruolo_tirocini}, Ergon Informatica considera il tirocinio uno degli strumenti principali per l'esplorazione tecnologica. Per alimentare costantemente questo ecosistema di sperimentazione, l'organizzazione ha strutturato una solida rete di collaborazioni con le principali istituzioni formative del territorio, accogliendo annualmente oltre una decina di tirocinanti. 

\subsection{Collaborazioni con le istituzioni formative}
\label{subsec:collaborazioni_formative}

L'azienda seleziona candidati da provenienze volutamente eterogenee: esse spaziano dai percorsi accademici tradizionali, come l'Università degli Studi di Padova e l'Università Ca' Foscari di Venezia, fino a iter maggiormente professionalizzanti come i vari Istituti Tecnici Superiori (ITS) della regione. 

Questa diversificazione consente all'azienda di individuare con continuità nuovi talenti e, ancor più importante, le permette di importare approcci mentali e operativi differenti. Come illustro nella Tabella~\ref{tab:confronto_profili_tirocinio}, mentre i profili universitari tendono ad apportare un maggiore rigore ingegneristico e metodologico, i profili tecnici introducono spesso una spiccata propensione per l'operatività immediata. 

\begin{table}[H]
    \centering
    \begin{tabular}{p{0.35\textwidth} p{0.55\textwidth}}
        \hline
        \textbf{Profilo del Tirocinante} & \textbf{Contributo Metodologico e Operativo} \\
        \hline
        Percorsi Accademici Universitari & Rigore ingegneristico, approfondimento teorico e approccio strutturato \\
        Istituti Tecnici Superiori & Immediatezza esecutiva, pragmatismo operativo e orientamento pratico \\
        \hline
    \end{tabular}
    \caption{Diversificazione degli approcci metodologici apportati dai tirocinanti.}
    \vspace{0.1cm}
    {\footnotesize Elaborazione originale dell'autore\par}
    \label{tab:confronto_profili_tirocinio}
\end{table}

L'inserimento di queste risorse consente inoltre all'organizzazione di osservare indirettamente l'evoluzione del panorama tecnologico e formativo esterno. Il confronto con studenti provenienti da percorsi differenti permette ai responsabili di raccogliere indicazioni sulle tecnologie maggiormente diffuse nei percorsi di studio, sugli strumenti emergenti e sugli approcci metodologici più recenti. In tale prospettiva, i tirocinanti diventano il collegamento tra il contesto accademico e quello industriale.

\subsection{Integrazione del tirocinante nel contesto aziendale}
\label{subsec:integrazione_tirocinante}

Sulla base delle informazioni raccolte, ho potuto comprendere come Ergon progetti generalmente le attività di tirocinio evitando mansioni esclusivamente esecutive o a basso valore formativo. I responsabili concepiscono le iniziative per coniugare le esigenze concrete dell'organizzazione con gli interessi e il percorso del candidato, con l'obiettivo di generare un coinvolgimento attivo e favorire la validazione di soluzioni con potenziale impatto reale. 

Questo approccio favorisce uno scambio di competenze all'interno dell'azienda: i tirocinanti beneficiano della profonda conoscenza del dominio applicativo e delle logiche commerciali trasferite dai colleghi esperti, mentre le figure \textit{senior} entrano in contatto con tecnologie emergenti e paradigmi moderni portati dai più giovani. In questo scenario, il tirocinante non assume il ruolo di semplice esecutore, ma contribuisce attivamente allo studio di strumenti e approcci alternativi, offrendo all'organizzazione l'opportunità di osservare nuove modalità operative e valutarne l'eventuale adozione futura.

\section{Ciclo di vita del tirocinio}
\label{sec:ciclo_vita_tirocinio}
Affinché l'apporto innovativo dei tirocinanti si traduca in un \textit{asset} tangibile e non in una dispersione di energie, Ergon ha istituzionalizzato un modello operativo rigoroso: l'azienda inserisce ogni progetto formativo in un ciclo di vita continuo, la cui architettura prevede un chiaro "prima" e un altrettanto definito "dopo" rispetto all'orizzonte temporale dello studente in sede.

\subsection{La genesi del progetto}
\label{subsec:genesi_necessita}

Il ciclo di vita di un progetto di \textit{stage} inizia molto prima dell'arrivo del candidato. La direzione definisce le iniziative partendo da problematiche operative concrete, che scaturiscono da due fonti primarie: 
\begin{itemize}
    \item \textbf{Impulsi di mercato:} specifiche richieste della clientela per funzionalità moderne che il gestionale storico non è in grado di erogare agilmente, come ad esempio cruscotti \textit{web} interattivi, interfacce \textit{mobile} per gli operatori logistici o portali di comunicazione \gls{b2b} --- ovvero sistemi concepiti per l'interazione commerciale diretta tra aziende, senza l'intermediazione del consumatore finale.
    \item \textbf{Esigenze infrastrutturali interne:} colli di bottiglia architettonici che la dirigenza tecnica rileva e che richiedono un'indagine preliminare su strumenti più efficienti --- quali moderni \gls{orm}, una tecnica di programmazione che mappa le tabelle di una base di dati relazionale su oggetti del codice sorgente per semplificarne la persistenza, o architetture \gls{restful}, uno stile architetturale per lo sviluppo di servizi \textit{web} standardizzati e disaccoppiati basati sul protocollo HTTP --- prima di una loro adozione definitiva.
\end{itemize}

Quando emerge una di queste esigenze, il \textit{team} tecnico ne delimita il perimetro funzionale e individua uno \textit{stack} tecnologico coerente con gli obiettivi esplorativi del progetto. Il gruppo di lavoro traduce così il problema operativo in un'attività di tirocinio delimitata e verificabile, destinata a produrre sia un risultato concreto sia elementi utili alla valutazione di possibili evoluzioni tecnologiche.

\subsection{Il processo di validazione tecnica}
\label{subsec:validazione_tecnica}

Secondo il modello organizzativo illustrato dal tutor aziendale, la conclusione formale del tirocinio, coincidente con il completamento del monte ore previsto, rappresenta generalmente per l'azienda l'inizio di una fase di consolidamento dei risultati ottenuti. In questa fase, l'azienda persegue l'obiettivo di produrre un modulo sufficientemente autonomo da poter essere valutato in termini tecnici e organizzativi; anche nei casi in cui il livello di completezza non consenta un impiego immediato, la dirigenza conserva e considera normalmente il lavoro svolto come base di riferimento per sviluppi successivi.

Al termine dello \textit{stage}, gli sviluppatori \textit{senior} sottopongono il codice sorgente, le scelte architetturali effettuate, la documentazione redatta e le pratiche di collaudo a una rigorosa procedura di \textit{code review}. I responsabili tecnici non si limitano a verificare il funzionamento dell'interfaccia, ma effettuano un \textit{audit} approfondito per valutare la manutenibilità del codice, la sicurezza delle \gls{api} --- le interfacce di programmazione che permettono ad applicativi diversi di comunicare e scambiare dati in modo standardizzato --- e la reale scalabilità delle tecnologie proposte dal tirocinante.

\subsection{Integrazione e sviluppo iterativo}
\label{subsec:integrazione_sviluppo}

A seguito del processo di validazione, il codice generato confluisce nel patrimonio informativo aziendale seguendo una di due possibili traiettorie di integrazione, a seconda della maturità del prodotto e della complessità del dominio affrontato:

\begin{enumerate}
    \item \textbf{Messa in produzione diretta:} Se l'applicativo soddisfa i rigidi requisiti di affidabilità aziendali e lo \textit{stack} tecnologico impiegato ha dimostrato la robustezza attesa, il \textit{software} entra in una fase di rifinitura interna. Gli sviluppatori \textit{senior} collegano il modulo sperimentale ai sistemi di persistenza di produzione e procedono con il rilascio progressivo ai clienti. In questo scenario, il tirocinio ha fornito all'azienda un prodotto commerciale "chiavi in mano".
    \item \textbf{Sviluppo iterativo e passaggio di consegne:} Qualora il progetto affronti un dominio applicativo troppo vasto per essere esaurito nel tempo a disposizione, o se la tecnologia investigata richiede ulteriori validazioni architetturali, il lavoro dello stagista assume il ruolo di \textit{baseline} (versione di riferimento). La dirigenza tecnica congela tale base di partenza e la riassegna successivamente come punto di inizio per un tirocinante futuro. Questo approccio a "staffetta" permette di aggredire problemi complessi suddividendoli in più cicli formativi, generando iterativamente valore incrementale.
\end{enumerate}

In entrambi gli scenari, il tirocinio contribuisce ad arricchire il patrimonio tecnico e organizzativo dell'azienda: da un lato attraverso la produzione di componenti o conoscenze riutilizzabili, dall'altro fornendo elementi utili per orientare le future decisioni di evoluzione tecnologica del sistema informativo. La Figura~\ref{fig:ciclo_tirocini_ergon} riassume visivamente l'intero ciclo di vita esplorativo, dalla raccolta delle necessità iniziali fino alla diramazione conclusiva tra produzione e sviluppo iterativo.

\begin{figure}[H]
    \centering
    \resizebox{\textwidth}{!}{
    \begin{tikzpicture}[
        node distance=1.5cm and 2cm,
        box/.style={rectangle, draw=black, thick, fill=blue!5, text width=4cm, align=center, minimum height=1.2cm, font=\small},
        decision/.style={diamond, aspect=2, draw=black, thick, fill=orange!10, text width=2.5cm, align=center, font=\footnotesize, inner sep=0pt},
        arrow/.style={thick, ->, >=stealth}
    ]
        % Nodi
        \node[box] (input) {\textbf{Genesi}\\Impulsi mercato o\\esigenze interne};
        \node[box, below=1cm of input] (sviluppo) {\textbf{Fase di Tirocinio}\\Sviluppo PoC / MVP};
        \node[box, below=1cm of sviluppo] (review) {\textbf{Validazione}\\\textit{Code Review} e \textit{Audit}};
        \node[decision, below=1cm of review] (bivio) {Maturità\\Progetto};
        \node[box, left=1cm of bivio] (prod) {\textbf{Messa in Produzione}\\Integrazione e rilascio};
        \node[box, right=1cm of bivio] (iter) {\textbf{Sviluppo Iterativo}\\\textit{Baseline} per futuri \textit{stage}};

        % Frecce
        \draw[arrow] (input) -- (sviluppo);
        \draw[arrow] (sviluppo) -- (review);
        \draw[arrow] (review) -- (bivio);
        \draw[arrow] (bivio) -- node[above=0.15cm, font=\footnotesize] {} (prod);
        \draw[arrow] (bivio) -- node[above=0.15cm, font=\footnotesize] {} (iter);
        
        % Freccia di ritorno per sviluppo iterativo
        \draw[arrow, dashed] (iter.north) |- node[above right=0.15cm, font=\footnotesize] {Nuovo ciclo} (sviluppo.east);
    \end{tikzpicture}
    }
    \caption{Rappresentazione schematica del ciclo di vita dei progetti in Ergon.}
    \vspace{0.1cm}
    {\footnotesize Elaborazione originale dell'autore\par}
    \label{fig:ciclo_tirocini_ergon}
\end{figure}

\section{Il dominio applicativo del progetto Kargo}
\label{sec:dominio_logistico}

In aderenza alle logiche di assegnazione precedentemente descritte, il mio progetto di tirocinio --- denominato "Kargo" --- trae origine da una specifica esigenza di mercato. Un cliente di Ergon Informatica, operante nel settore della distribuzione di bevande, ha manifestato la necessità di risolvere una forte criticità legata alla logistica \textit{inbound}: la congestione delle banchine di scarico. 

Nella gestione quotidiana dei magazzini, l'arrivo non programmato dei vettori e la simultaneità delle consegne generano inefficienze a catena. I magazzinieri subiscono picchi di lavoro imprevedibili, con conseguenti ritardi nello stoccaggio, mentre i corrieri sono costretti a lunghi periodi di attesa nei piazzali prima di poter effettuare le operazioni di scarico, con un inevitabile aggravio dei costi operativi.

L'azienda, incaricata di fornire una soluzione, aveva inizialmente pianificato uno sviluppo interno tradizionale. Tuttavia, riconoscendo in questa specifica problematica le dimensioni ideali per uno \textit{stage} e dopo aver valutato la necessità di esplorare nuove tecnologie \textit{web}, la dirigenza ha deciso di convertire la commessa in un progetto di tirocinio esplorativo.

\subsection{Obiettivi e vincoli del progetto Kargo}
\label{subsec:obiettivi_vincoli_kargo}

Attraverso un confronto diretto con il mio tutor aziendale, abbiamo definito congiuntamente il perimetro operativo del progetto Kargo, tracciando obiettivi funzionali chiari e quantificabili, e imponendo vincoli architetturali strategici.

\textbf{L'obiettivo funzionale aziendale} consisteva nello sviluppo di un'applicazione \textit{web} in grado di digitalizzare e automatizzare la prenotazione degli \textit{slot} temporali per lo scarico merci. L'applicativo doveva consentire ai corrieri di monitorare le disponibilità delle banchine in tempo reale, prenotare un appuntamento in modo sicuro e fornire all'amministrazione uno strumento di pianificazione ordinata. Per valutare oggettivamente il raggiungimento del risultato, Ergon ha fissato l'aspettativa di ricevere, entro le 300 ore previste dal tirocinio, un \gls{mvp} --- ovvero una versione iniziale del \textit{software} limitata alle funzionalità essenziali per completare un intero ciclo logistico (dalla registrazione del corriere alla conferma dello \textit{slot}) e raccogliere i riscontri dei committenti --- strutturalmente idoneo a essere rifinito e presentato al cliente.

\textbf{I vincoli tecnologici e architetturali} stabiliti da Ergon sono stati altrettanto rigorosi. L'azienda ha richiesto espressamente l'utilizzo del \textit{framework} React per lo sviluppo del livello di presentazione (\textit{frontend}), una tecnologia inedita rispetto agli standard interni. Inoltre, ha imposto un vincolo di stretta modularità architetturale: ho dovuto progettare il sistema per renderlo agilmente integrabile, in un secondo momento, con il gestionale \textit{legacy} proprietario descritto in Sezione~\ref{subsec:ecosistema_applicativo}. Questo requisito ha reso necessario un approccio disaccoppiato, basato sullo sviluppo di API in ambiente .NET per il livello logico. La Figura~\ref{fig:vincoli_architettura_kargo} illustra questa netta separazione strutturale, evidenziando il perimetro operativo assegnato al progetto Kargo e le relative modalità di comunicazione con il preesistente livello di persistenza.

\begin{figure}[H]
    \centering
    \resizebox{\textwidth}{!}{
    \begin{tikzpicture}[
        node distance=3cm,
        box/.style={rectangle, draw=black, thick, text width=4cm, align=center, minimum height=1.8cm, font=\small},
        arrow/.style={thick, <->, >=stealth}
    ]
        % Nodi con colori distintivi leggeri
        \node[box, fill=cyan!10] (frontend) {\textbf{Livello Presentazione}\\Applicazione Web\\(React)};
        \node[box, fill=green!10, right=2.5cm of frontend] (api) {\textbf{Livello Logico}\\Servizi RESTful\\(.NET C\#)};
        \node[box, fill=red!10, right=2.5cm of api] (legacy) {\textbf{Livello Persistenza}\\\textit{Database} Informix};

        % Frecce
        \draw[arrow] (frontend) -- node[above=0.15cm, font=\footnotesize] {Richieste HTTP} node[below=0.15cm, font=\footnotesize] {JSON} (api);
        \draw[arrow] (api) -- node[above=0.15cm, font=\footnotesize] {Integrazione} node[below=0.15cm, font=\footnotesize] {Legacy} (legacy);
        
        % Riquadro di raggruppamento ancorato automaticamente
        \draw[dashed, thick, draw=gray] 
            ([xshift=-0.5cm, yshift=0.9cm]frontend.north west) 
            rectangle 
            ([xshift=0.5cm, yshift=-0.5cm]api.south east);
            
        % Testo del perimetro
        \node[font=\footnotesize, text=gray, anchor=north west] 
            at ([xshift=-0.4cm, yshift=0.8cm]frontend.north west) 
            {Perimetro di sviluppo del progetto Kargo};

    \end{tikzpicture}
    }
    \caption{Vincoli di modularità architetturale imposti per il progetto Kargo.}
    \vspace{0.1cm}
    {\footnotesize Elaborazione originale dell'autore\par}
    \label{fig:vincoli_architettura_kargo}
\end{figure}

\section{Motivazioni e traguardi personali}
\label{sec:motivazioni_scelta}

Durante la fase di selezione delle proposte di \textit{stage}, la mia candidatura iniziale era in realtà orientata verso un differente progetto interno, relativo allo sviluppo di un sistema per il controllo degli accessi aziendali basato su tecnologia .NET MAUI. Tuttavia, a seguito di un colloquio orientativo con la direzione delle Risorse Umane, l'azienda mi ha proposto il progetto Kargo, illustrandone lo \textit{stack} tecnologico, le problematiche da risolvere e gli obiettivi. Ho quindi scelto di orientare il mio percorso formativo verso quest'ultima proposta per motivazioni di natura sia tecnologica sia metodologica.

In primo luogo, dal punto di vista tecnico, Kargo rappresentava l'unica iniziativa a proporre lo sviluppo di interfacce tramite React. Rispetto alle alternative basate su \textit{framework} maggiormente legati all'ecosistema \textit{desktop} o \textit{mobile} ibrido, l'opportunità di confrontarmi con lo sviluppo \textit{web full-stack} moderno ha costituito per me un incentivo professionale determinante. 

In secondo luogo, dal punto di vista metodologico, la natura esplorativa del progetto si allineava perfettamente con i miei obiettivi accademici. Durante il percorso universitario, ed in particolare nel corso di Ingegneria del Software, ho sviluppato un forte interesse per le tematiche relative al \gls{system_design} --- l'insieme delle metodologie di progettazione dell'architettura e dei componenti di sistemi informatici complessi --- e all'architettura del \textit{software}. Un progetto nato al di fuori dei rigidi binari della produzione standard mi offriva un maggiore margine di manovra: garantiva la libertà necessaria per progettare un sistema \textit{ex novo}, permettendomi di applicare con rigore i \textit{pattern} architetturali appresi e di analizzare criticamente le mie scelte ingegneristiche in un ambiente protetto.

A fronte di queste premesse, ho stabilito dei traguardi personali oggettivi, articolati nei seguenti punti:
\begin{itemize}
    \item \textbf{Padronanza tecnologica e metodologica:} raggiungere la piena autonomia operativa nello sviluppo di applicazioni basate su C\# e React, comprovata dalla consegna dell'MVP nei tempi stabiliti, e sviluppare parallelamente un metodo di studio efficiente per assimilare rapidamente nuovi linguaggi di programmazione.
    \item \textbf{Architettura e \textit{problem solving}:} progettare un'architettura \textit{software} applicando consapevolmente i principi di \textit{System Design}, che ottenga la validazione positiva del \textit{team} tecnico durante le procedure di \textit{code review}.
    \item \textbf{Gestione del ciclo di vita del \textit{software}:} condurre operativamente tutte le principali fasi di produzione, tracciando il passaggio dalla progettazione iniziale alla validazione finale.
\end{itemize}

Mi sono posto l'obiettivo di verificare quanto le conoscenze teoriche acquisite durante il corso di laurea fossero effettivamente applicabili in un contesto industriale, consolidandole attraverso lo svolgimento diretto delle attività progettuali. I traguardi definiti in questa sezione costituiranno quindi il riferimento metrico con cui, nel Capitolo~\ref{chap:conclusioni}, valuterò oggettivamente i risultati ottenuti al termine del tirocinio.